#include "main.typ"
\section{Một số bài tập}
% ================= CÂU 1 =================
\noindent {\color{myred}\textbf{Câu 1.}} Đường cong $y = \dfrac{1}{4}x^4 - (3m+1)x^2 + 2m + 2$ có ba điểm cực trị tạo thành tam giác $ABC$ sao cho tam giác $ABC$ nhận gốc tọa độ làm trọng tâm. Giá trị tham số $m$ gần nhất với giá trị nào?
\begin{tabbing}
A. $0,35$ \qquad\qquad \= B. $0,79$ \qquad\qquad \= C. $0,96$ \qquad\qquad \= D. $1,52$
\end{tabbing}

\vspace{0.3cm}

% ================= CÂU 2 =================
\noindent {\color{myred}\textbf{Câu 2.}} Ký hiệu $S$ là tập hợp tất cả các giá trị $m$ để đường cong $y = x^4 - 2mx^2 + 3$ có ba điểm cực trị tạo thành tam giác $ABC$ sao cho tam giác tồn tại một góc $30^\circ$. Tổng tất cả các phần tử của $S$ gần nhất với giá trị nào?
\begin{tabbing}
A. $3$ \qquad\qquad\qquad \= B. $4$ \qquad\qquad\qquad \= C. $6$ \qquad\qquad\qquad \= D. $7$
\end{tabbing}

\vspace{0.3cm}

% ================= CÂU 3 =================
\noindent {\color{myred}\textbf{Câu 3.}} Đường cong $y = x^4 - 2mx^2 - m - 1$ có ba điểm cực trị lập thành một tam giác $ABC$ sao cho tam giác tồn tại một góc $120^\circ$. Giá trị tham số $m$ là
\begin{tabbing}
A. $m = -\dfrac{1}{\sqrt[3]{3}}$ \qquad \= B. $m = \sqrt[3]{3}$ \qquad\qquad \= C. $m = -\dfrac{1}{\sqrt[3]{2}}$ \qquad \= D. $m = -\dfrac{1}{\sqrt[3]{5}}$
\end{tabbing}

\vspace{0.3cm}

% ================= CÂU 4 =================
\noindent {\color{myred}\textbf{Câu 4.}} Đường cong $y = x^4 - 2mx^2 + 3$ có ba điểm cực trị lập thành tam giác $ABC$ cân sao cho độ dài cạnh bên gấp đôi độ dài cạnh đáy. Giá trị tham số $m$ là
\begin{tabbing}
A. $m = \sqrt[3]{15}$ \qquad\quad \= B. $m = \sqrt[3]{3}$ \qquad\qquad \= C. $m = -\dfrac{1}{\sqrt[3]{2}}$ \qquad \= D. $m = \sqrt[3]{10}$
\end{tabbing}

\vspace{0.3cm}

% ================= CÂU 5 =================
\noindent {\color{myred}\textbf{Câu 5.}} Đường cong $y = x^4 - 2mx^2 + 3$ có ba điểm cực trị lập thành tam giác $ABC$ cân sao cho độ dài đường cao xuất phát từ đỉnh bằng độ dài cạnh đáy. Giá trị tham số $m$ là
\begin{tabbing}
A. $m = \sqrt[3]{15}$ \qquad\quad \= B. $m = \sqrt[3]{3}$ \qquad\qquad \= C. $m = \sqrt[3]{4}$ \qquad\quad \= D. $m = \sqrt[3]{10}$
\end{tabbing}

\vspace{0.3cm}

% ================= CÂU 6 =================
\noindent {\color{myred}\textbf{Câu 6.}} Tồn tại duy nhất giá trị $m = k$ để đường cong $y = x^4 - 2(m^2 - m + 1)x^2 + m - 1$ có ba điểm cực trị phân biệt sao cho khoảng cách giữa hai điểm cực tiểu đạt giá trị nhỏ nhất. Mệnh đề nào sau đây đúng?
\begin{tabbing}
A. $0 < k < 1$ \qquad\quad \= B. $2 < k < 3$ \qquad\quad \= C. $k > 4$ \qquad\qquad \= D. $k = 3$
\end{tabbing}
% ================= CÂU 7 =================
\noindent {\color{myred}\textbf{Câu 7.}} Tìm điều kiện của tham số $m$ để hàm số $y = (1-m)x^4 + 2(m+3)x^2 + 1$ có đúng một điểm cực tiểu và không có điểm cực đại.
\begin{tabbing}
A. $m < 1$ \qquad\qquad\qquad \= B. $m < -3$ \qquad\qquad\qquad \= C. $m > 1$ \qquad\qquad\qquad \= D. $-3 \le m \le 1$
\end{tabbing}

\vspace{0.3cm}

% ================= CÂU 8 =================
\noindent {\color{myred}\textbf{Câu 8.}} Hàm số $y = \dfrac{m+1}{2}x^4 - mx^2 + \dfrac{5}{2}$ có cực tiểu và không có cực đại khi $m$ thuộc khoảng $(a;b)$. Tìm $b-a$.
\begin{tabbing}
A. $1$ \qquad\qquad\qquad\qquad \= B. $2$ \qquad\qquad\qquad\qquad \= C. $3$ \qquad\qquad\qquad\qquad \= D. $4$
\end{tabbing}

\vspace{0.3cm}

% ================= CÂU 9 =================
\noindent {\color{myred}\textbf{Câu 9.}} Hàm số $y = mx^4 + (m-1)x^2 + 1 - 2m$ chỉ có một cực trị khi $m \le a$ hoặc $m \ge b$. Tính giá trị $M = a + b$.
\begin{tabbing}
A. $1$ \qquad\qquad\qquad\qquad \= B. $2$ \qquad\qquad\qquad\qquad \= C. $3$ \qquad\qquad\qquad\qquad \= D. $4$
\end{tabbing}

\vspace{0.3cm}

% ================= CÂU 10 =================
\noindent {\color{myred}\textbf{Câu 10.}} Tìm điều kiện của $m$ để hàm số $y = \dfrac{1}{2}x^4 - mx^2 + \dfrac{3}{2}$ chỉ có cực tiểu và không có cực đại.
\begin{tabbing}
A. $m \le 0$ \qquad\qquad\qquad \= B. $m > 0$ \qquad\qquad\qquad \= C. $m < 2$ \qquad\qquad\qquad \= D. $m < 0$
\end{tabbing}

\vspace{0.3cm}

% ================= CÂU 11 =================
\noindent {\color{myred}\textbf{Câu 11.}} Tìm giá trị lớn nhất của $m$ để hàm số $y = -x^4 + 2(m+2)x^2 - 2m - 3$ chỉ có cực đại, không có cực tiểu.
\begin{tabbing}
A. $m = -2$ \qquad\qquad\quad \= B. $m = -1$ \qquad\qquad\quad \= C. $m = 0$ \qquad\qquad\qquad \= D. $m = 3$
\end{tabbing}

\vspace{0.3cm}

% ================= CÂU 12 =================
\noindent {\color{myred}\textbf{Câu 12.}} Trong khoảng $(-5; 4)$ tồn tại bao nhiêu giá trị nguyên của $m$ để hàm số $y = -x^4 + 2(m+2)x^2 - 2m - 3$ chỉ có cực đại, không có cực tiểu?
\begin{tabbing}
A. $3$ giá trị. \qquad\qquad\quad \= B. $2$ giá trị. \qquad\qquad\quad \= C. $1$ giá trị. \qquad\qquad\quad \= D. $4$ giá trị.
\end{tabbing}

\vspace{0.3cm}

% ================= CÂU 13 =================
\noindent {\color{myred}\textbf{Câu 13.}} Hàm số $y = x^4 + (m+3)x^3 + 2(m+1)x^2$ có cực đại tại $x = k$ và $k \le 0$. Trong khoảng $[-2017; 2017]$ tồn tại bao nhiêu giá trị nguyên của $k$ thỏa mãn yêu cầu bài toán?
\begin{tabbing}
A. $4034$ \qquad\qquad\qquad \= B. $4033$ \qquad\qquad\qquad \= C. $4032$ \qquad\qquad\qquad \= D. $4035$
\end{tabbing}

\vspace{0.3cm}

% ================= CÂU 14 =================
\noindent {\color{myred}\textbf{Câu 14.}} Hàm số $y = mx^4 - (2m+1)x^2 + 1$ có một điểm cực đại khi $m \ge k$. Giá trị $k$ nằm trong khoảng nào?
\begin{tabbing}
A. $-1 < k < 0$ \qquad\qquad \= B. $0 < k < 1$ \qquad\qquad \= C. $2 < k < 4$ \qquad\qquad \= D. $-3 < k < -1$
\end{tabbing}

\vspace{0.3cm}

% ================= CÂU 15 =================
\noindent {\color{myred}\textbf{Câu 15.}} Tìm giá trị lớn nhất của tham số $m$ để đường cong $y = \dfrac{1}{6}(m+2)x^4 - (m-1)x^2 + 5$ có đúng một cực tiểu.
\begin{tabbing}
A. $m = 1$ \qquad\qquad\qquad \= B. $m = 0$ \qquad\qquad\qquad \= C. $m = 2$ \qquad\qquad\qquad \= D. $m = 0,5$
\end{tabbing}
\noindent {\color{myred}\textbf{Câu 15.}} Cho hàm số $y=x^4-2mx^2+3m-2$ (với $m$ là tham số). Có bao nhiêu giá trị của tham số $m$ để đồ thị hàm số có ba điểm cực trị đều nằm trên các trục tọa độ?
\begin{tabbing}
A. $2.$ \qquad\qquad\qquad \= B. $0.$ \qquad\qquad\qquad \= C. $3.$ \qquad\qquad\qquad \= D. $1.$
\end{tabbing}

\vspace{0.3cm}

\noindent {\color{myred}\textbf{Câu 16.}} Biết $m=m_0; m_0 \in \mathbb{R}$ là giá trị của tham số $m$ để đồ thị hàm số $y=x^4+2mx^2+1$ có ba điểm cực trị tạo thành một tam giác vuông. Khẳng định nào sau đây đúng?
\begin{tabbing}
A. $m_0 \in (0;3).$ \qquad\qquad\qquad \= B. $m_0 \in [-5;-3).$ \qquad\qquad\qquad \= C. $m_0 \in (-3;0].$ \qquad\qquad\qquad \= D. $m_0 \in (3;7).$
\end{tabbing}

\vspace{0.3cm}

\noindent {\color{myred}\textbf{Câu 17.}} Cho hàm số $y=x^4-2(m^2+m+1)x^2+m$ có đồ thị $(C)$. Tìm $m$ để đồ thị hàm số $(C)$ có $3$ điểm cực trị và khoảng cách giữa hai điểm cực tiểu nhỏ nhất.
\begin{tabbing}
A. $m=\dfrac{1}{2}.$ \qquad\qquad\qquad \= B. $m=-\dfrac{1}{2}.$ \qquad\qquad\qquad \= C. $m=\sqrt{3}.$ \qquad\qquad\qquad \= D. $m=0.$
\end{tabbing}

\vspace{0.3cm}

\noindent {\color{myred}\textbf{Câu 18.}} Để đồ thị hàm số $y=x^4-2mx^2+m-1$ có ba điểm cực trị tạo thành một tam giác có diện tích bằng $2$, giá trị của tham số $m$ thuộc khoảng nào sau đây?
\begin{tabbing}
A. $(2;3).$ \qquad\qquad\qquad \= B. $(-1;0).$ \qquad\qquad\qquad \= C. $(0;1).$ \qquad\qquad\qquad \= D. $(1;2).$
\end{tabbing}

\vspace{0.3cm}

\noindent {\color{myred}\textbf{Câu 19.}} Cho hàm số $f(x)=x^4-2mx^2+4-2m^2$. Có tất cả bao nhiêu số nguyên $m \in (-10;10)$ để hàm số $y=|f(x)|$ có đúng $3$ điểm cực trị?
\begin{tabbing}
A. $6.$ \qquad\qquad\qquad \= B. $8.$ \qquad\qquad\qquad \= C. $9.$ \qquad\qquad\qquad \= D. $7.$
\end{tabbing}

\vspace{0.3cm}
\vfill
